\documentclass[11pt,a4paper]{article}
\usepackage[utf8]{inputenc}
\usepackage[T1]{fontenc}
\usepackage{amsmath,amssymb}

\title{{{title}}}
\author{Lab Partner 1 \and Lab Partner 2\\Department of Physics}
\date{\today}

\begin{document}
\maketitle

\section{Objective}
The purpose of this laboratory experiment is to quantitatively measure thermal dissipation in driven systems and verify classical thermodynamic bounds as discussed in the literature~\cite{knuth1984}.

\section{Theoretical Background}
The governing equation relates entropy change to heat transfer:
\begin{equation}
\Delta S \ge \int_0^t \frac{\delta Q}{T}
\end{equation}

\section{Apparatus and Procedure}
\begin{enumerate}
  \item Calibrate thermal sensors against standard reference reservoirs.
  \item Initialize the sample under controlled ambient pressure.
  \item Record sensor trajectories at 1~kHz sampling rate.
\end{enumerate}

\section{Data and Results}
Measurements collected over multiple trials demonstrated high consistency with theoretical limits. Uncertainty bounds were evaluated at the $95\%$ confidence interval.

\section{Discussion and Error Analysis}
Systematic uncertainties were dominated by thermal sensor response times, while random noise remained well within nominal operating tolerance.

\section{Conclusion}
The experimental results successfully validate the theoretical prediction within the established error margins.

\bibliographystyle{plain}
\bibliography{references}

\end{document}
